% SOURCE_STATEMENT: "Consider the Korteweg--de Vries (KdV) equation u_t + u u_x + u_{xxx} =0 with initial data u(x, t=0) = u_0(x),\\qquad x\\in \\mathbb{R}. In the 1970's, McKean and Trubowitz proved the remarkable result that if the initial data $u_0$ is periodic, $u_0(x+p)= u_0(x)$ for some $p>0$, then the solution $u(x,t)$ of (eq1) is almost periodic in time. This result leads to the following natural conjecture: The same is true if $u_0(x)$ is almost periodic, i.e., if the initial data is almost periodic in space, the solution evolves almost periodically in time. This is a very challenging problem with many new and unique features. For example, for $u_0(x)$ almost periodic, one does not know, by standard PDE methods, whether a solution $u(x,t)$ exists, even for a small time, never mind globally. In spectacular and seminal work over the last few years, D. Damanik and M. Goldstein, joined later by I. Binder and M. Lukic, have partially resolved this conjecture in the affirmative, namely, for small quasi-periodic analytic initial data with Diophantine frequencies, unique global solutions exist and are almost periodic in time. The key insight of Damanik et al.\\ was how to utilize the formal integrability of KdV in an effective analytical manner. Not only are the arguments in Damanik et al.\\ applicable to other (formally) integrable systems (for example proves a Toda lattice analog of the KdV theorem), but some of their arguments are also applicable to non-integrable systems. It is a natural, open and very challenging conjecture that the result of Damanik et al.\\ remains true for (suitable) perturbations of KdV, i.e., for (suitable) perturbations of KdV, the Cauchy problem with (suitable) almost periodic initial data $u_0(x)$, has a unique global solution $u(x,t)$, which evolves almost periodically in time. After all, the finite-dimensional analog of this result is precisely the content of classical KAM theory. We note that such results for perturbations of KdV and the nonlinear Schr\\\"{o}dinger equation (NLS) with periodic initial data $u_0(x)$, have already been proved by S. Kuksin, T. Kappeler and others. With the discovery of quasi-crystals, one can anticipate, in particular, that interest in PDE problems on quasi-periodic backgrounds, will surely grow."
% Preparation and final self-audit: CONTRIBUTING.md prompts C and D; source/audit artifacts retained in the local physics_statements_2026_10_11 batch.
% !TeX program = lualatex
\documentclass[11pt]{article}
\usepackage[margin=25mm]{geometry}
\usepackage{fontspec}
\setmainfont{Latin Modern Roman}
\usepackage{amsmath,amssymb,mathtools}
\usepackage{unicode-math}
\setmathfont{Latin Modern Math}
\usepackage{xurl,hyperref}
\hypersetup{colorlinks=true,urlcolor=blue}
\setcounter{secnumdepth}{0}
\setlength{\parindent}{0pt}
\setlength{\parskip}{5pt}
\setlength{\emergencystretch}{3em}
\newcommand{\sref}[2]{\hyperlink{source:#1:#2}{[S#2]}}
\newcommand{\cataloguescope}[1]{\paragraph{Recorded target.}#1}
\newcommand{\R}{\mathbb R}
\newcommand{\C}{\mathbb C}
\newcommand{\N}{\mathbb N}
\newcommand{\Z}{\mathbb Z}
\newcommand{\Q}{\mathbb Q}
\newcommand{\F}{\mathbb F}
\newcommand{\eps}{\varepsilon}
\begin{document}
\section{KdV with almost periodic initial data}\label{statement}

\subsection{Definitions and mathematical statement}

\paragraph{Shared notation.}$\mathbb N=\{1,2,\ldots\}$; the standard number fields have their usual meanings. All additional conventions are stated locally.


\paragraph{Definitions and notation.}A bounded function is Bohr almost periodic if its translates form a precompact family in the uniform norm; equivalently it is uniformly approximable by trigonometric polynomials. Quasiperiodic data depend on finitely many torus frequencies. The KdV normalization used here is $u_t+uu_x+u_{xxx}=0$. Almost periodic evolution in time requires a specified function-space topology, not merely boundedness of a solution.

\paragraph{Catalogue wording retained for traceability.}
Consider the Korteweg--de Vries (KdV) equation \[u_t+uu_x+u_{xxx}=0,\qquad u(x,0)=u_0(x),\qquad x\in\mathbb R.\] In the 1970's, McKean and Trubowitz proved the remarkable result that if the initial data $u_0$ is periodic, $u_0(x+p)= u_0(x)$ for some $p>0$, then the solution $u(x,t)$ of this equation is almost periodic in time. This result leads to the following natural conjecture: The same is true if $u_0(x)$ is almost periodic, i.e., if the initial data is almost periodic in space, the solution evolves almost periodically in time. This is a very challenging problem with many new and unique features. For example, for $u_0(x)$ almost periodic, one does not know, by standard PDE methods, whether a solution $u(x,t)$ exists, even for a small time, never mind globally. In spectacular and seminal work over the last few years, D. Damanik and M. Goldstein, joined later by I. Binder and M. Lukic, have partially resolved this conjecture in the affirmative, namely, for small quasi-periodic analytic initial data with Diophantine frequencies, unique global solutions exist and are almost periodic in time. The key insight of Damanik et al.\ was how to utilize the formal integrability of KdV in an effective analytical manner. Not only are the arguments in Damanik et al.\ applicable to other (formally) integrable systems (for example proves a Toda lattice analog of the KdV theorem), but some of their arguments are also applicable to non-integrable systems. It is a natural, open and very challenging conjecture that the result of Damanik et al.\ remains true for (suitable) perturbations of KdV, i.e., for (suitable) perturbations of KdV, the Cauchy problem with (suitable) almost periodic initial data $u_0(x)$, has a unique global solution $u(x,t)$, which evolves almost periodically in time. After all, the finite-dimensional analog of this result is precisely the content of classical KAM theory. We note that such results for perturbations of KdV and the nonlinear Schr\"{o}dinger equation (NLS) with periodic initial data $u_0(x)$, have already been proved by S. Kuksin, T. Kappeler and others. With the discovery of quasi-crystals, one can anticipate, in particular, that interest in PDE problems on quasi-periodic backgrounds, will surely grow.
\paragraph{Scope and source qualification.}The original source records partial progress for small analytic quasiperiodic data with Diophantine frequencies. Its words suitable perturbations and suitable data are part of a research programme; no unsupported universal perturbation class is imposed.
\subsection{Short English statement}
Prove global and time-almost-periodic evolution for suitable almost-periodic KdV data and establish persistence under appropriate perturbations.
\subsection{Sources}
\begin{itemize}
\item[\hypertarget{source:statement:1}{S1}] ULAM.AI and the UnsolvedMath Contributors, UnsolvedMath, catalogue entry. Catalogue curation under CC BY 4.0; definitions and English exposition added locally.
\url{https://huggingface.co/datasets/ulamai/UnsolvedMath}.
\item[\hypertarget{source:statement:2}{S2}] Percy Deift, Some Open Problems in Random Matrix Theory and the Theory of Integrable Systems. II, SIGMA 13 (2017), 016, Problem 1.
\url{https://arxiv.org/abs/1703.04931}.
\item[\hypertarget{source:statement:3}{S3}] Primary full-text formulation and 2017 historical progress discussion.
\url{https://arxiv.org/html/1703.04931v1}.
\end{itemize}
\label{end:statement}
\end{document}
